\bookchapter{第三十二章：手机安全急救手册}

\chapterintro{这一章是一份“应急预案”。希望你永远用不上，但如果真的遇到手机丢失、被盗、被诈骗的情况，这份手册能帮你把损失降到最低；地震、暴雨等灾害来临时，手机同样是接收预警、定位求救、保持联络的工具。建议你现在就通读一遍，然后把关键步骤截图保存，或者打印出来放在家里——因为出事的时候你的手机可能已经不在手上了，灾害时则可能没电或没信号。}

\booksection{场景零：地震 / 极端天气——手机怎么帮你保命}

防的是人身安全，不是账号资金。预警设置详见第十三章「天气预警、自然灾害预警与地震预警」；这里只讲\textbf{警报响起之后}该怎么做。

\textbf{地震警报响了：先躲后撤，别只顾着拍视频}\par

\begin{enumerate}
  \item \textbf{就近避险}：趴下、掩护、抓牢（桌下、承重墙角、坚固家具旁），护住头颈；远离玻璃、外墙、吊灯。
  \item \textbf{不要盲目冲楼}：主震往往只有十几秒到几十秒，楼梯间可能更危险；电梯绝对不能用。
  \item \textbf{震停后再有序撤离}：走楼梯、注意余震与掉落物；到开阔地，远离电线杆、广告牌、危房。
  \item \textbf{手机用在刀刃上}：给紧急联系人发简短消息（「我在XX，安全/受伤，需要帮助」）比打长电话更省电；能发短信/微信时优先文字。
  \item \textbf{没信号时}：见第七章卫星 SOS / 卫星消息（机型支持时）；或到高处、开阔地再试；部分场景下无网蓝牙对讲亦有帮助。
  \item \textbf{关闭非必要功能}：应急时开飞行模式再定时开、关蓝牙扫描、降亮度，能多撑几小时；把最后一格电留给报平安与救援信息。
\end{enumerate}

\begin{tcbwarn}{暴雨 / 台风 / 山洪：预警阶段就要行动}
\begin{itemize}
  \item 看到\textbf{红色预警}：尽量减少外出，远离地下空间、下穿桥、河道、陡坡。
  \item 在户外被困：往坚固建筑物高处走，勿强行涉水（积水可能有井盖缺失与漏电）；打 119/110 时说清地标与水位。
  \item 出行前：用天气 App 与当地预警公众号确认风险；长途自驾避开山区临水临崖路段。
  \item 家人互报：约定「若失联，以最后微信/短信定位为准」，避免各自疯狂拨号耗光电量。
\end{itemize}

\end{tcbwarn}

\begin{tcbpractice}{平时就做好的 4 件事}
\begin{enumerate}
  \item 开启系统地震预警 / 正规预警 App，并允许其突破静音。
  \item 在通讯录置顶 2-3 位「紧急联系人」，锁屏界面可显示（iOS/安卓医疗急救卡或锁屏信息）。
  \item 家里和包里备充电宝；长途出行前充满电并带线。
  \item 记住至少一个「不用通讯录也能打」的号码（家人手机号），写在纸上放钱包。
\end{enumerate}

\end{tcbpractice}

\booksection{场景一：手机丢了/被偷了——黄金30分钟}

手机丢失后的前30分钟是“黄金时间”，处理得当可以保住你的钱财和隐私。按以下顺序操作，一步都不要跳：

\textbf{丢失后30分钟标准操作流程}\par

\begin{enumerate}
  \item \textbf{借别人手机拨打自己号码}（第1分钟）：可能只是掉在附近，好心人接了就能拿回来。
  \item \textbf{远程锁定手机}（第2-5分钟）：用别人的手机或电脑登录查找服务：
  \begin{itemize}
    \item iPhone：iCloud.com > 查找 > 选你的设备 > 标记为丢失
    \item 安卓各品牌：华为cloud.huawei.com / 小米i.mi.com / OPPO cloud.oppo.com 等
  \end{itemize}
  \item \textbf{挂失SIM卡}（第5-10分钟）：打运营商客服电话（移动10086/联通10010/电信10000），报手机丢失要求暂停SIM卡服务。这一步最关键——骗子拿到你的SIM卡就能接收验证码，重置你所有密码。
  \item \textbf{冻结支付账户}（第10-20分钟）：
  \begin{itemize}
    \item 支付宝：95188 挂失，或在别人手机上登录支付宝 > 我的 > 设置 > 安全中心 > 紧急冻结
    \item 微信：110 报警后联系微信客服，或在微信安全中心 > 冻结账号
    \item 银行App：拨打银行卡背面的客服电话，要求冻结网银
  \end{itemize}
  \item \textbf{通知亲友}（第20-25分钟）：在朋友圈或让家人帮忙发消息，通知亲友手机丢失，防止骗子用你的身份骗别人。
  \item \textbf{报警}（第25-30分钟）：拨打110或去最近的派出所报案，拿到报案回执。后续保险理赔、银行纠纷都需要这个回执。
\end{enumerate}

\booksection{场景二：被骗了——72小时黄金处理期}

发现自己被骗后，最关键是冷静。慌乱中做错误决定会导致二次损失。按以下步骤操作：

\begin{enumerate}
  \item \textbf{立即停止一切转账操作}：不管骗子说什么，不要再转一分钱。"再转一笔保证金就能退回之前的钱"是经典二次诈骗话术。
  \item \textbf{保留所有证据}：不要删聊天记录！不要删通话记录！截图保存聊天记录、转账记录、对方账号信息。这些都是报案和追回资金的关键证据。
  \item \textbf{立即拨打96110}：这是全国反诈专线。接线员会指导你做紧急止付——如果转账时间不长（通常24小时内），银行有可能冻结对方账户拦截资金。
  \item \textbf{拨打110报警}：带上手机、身份证、转账凭证去最近的派出所做笔录。
  \item \textbf{联系银行/支付平台}：告知被骗情况，要求冻结涉案账户。如果是信用卡被盗刷，要求银行启动拒付流程。
  \item \textbf{修改所有密码}：在另一台安全设备上修改你的Apple ID/品牌账号、微信、支付宝、银行App的密码，开启双重认证。
\end{enumerate}

\booksection{场景三：手机中了病毒/木马——怎么办？}

如果你发现手机出现以下症状——莫名其妙弹出广告、自动安装奇怪App、话费暴增、流量异常消耗——很可能是中了恶意软件：

\begin{itemize}
  \item \textbf{第一步：进入安全模式}——安卓手机长按电源键关机，再长按电源键开机，开机震动后长按音量减键不放，直到进入安全模式（左下角会显示"安全模式"字样）。安全模式下所有第三方App被禁用，方便排查。
  \item \textbf{第二步：卸载可疑App}——在安全模式下进入设置 > 应用管理，按时间排序，找到最近安装的不认识的App卸载掉。
  \item \textbf{第三步：安装安全App扫描}——下载"腾讯手机管家"或"360手机卫士"做一次全盘扫描。
  \item \textbf{第四步：恢复出厂设置}（终极方案）——如果以上方法无效，备份重要数据后恢复出厂设置。注意：恢复出厂会删除所有数据，一定要先备份照片、通讯录等重要文件。
  \item \textbf{iPhone基本不会中毒}——iOS的沙盒机制让恶意软件很难存活。如果你的iPhone出现异常，大概率是某个App的bug或设置问题，不是病毒。恢复出厂设置能解决99\%的问题。
\end{itemize}

\booksection{场景四：账号被盗——快速夺回控制权}

\begin{enumerate}
  \item \textbf{立即修改密码}：如果你还能登录账号，第一时间改密码。新密码要和旧密码完全不同。
  \item \textbf{踢掉其他设备}：在账号安全设置里找到"登录设备管理"，把不认识的设备全部踢下线。
  \item \textbf{检查登录记录}：查看最近7天的登录记录，如果有异地登录，说明账号确实被盗了。
  \item \textbf{开启双重认证}：如果你还没开双重认证，现在立刻开。这是防止账号再次被盗最有效的方法。
  \item \textbf{检查绑定信息}：检查账号绑定的手机号、邮箱有没有被改。如果被改了，用账号申诉功能找回。
  \item \textbf{检查支付设置}：如果被盗的是微信/支付宝，检查免密支付、自动扣款有没有被篡改，银行卡有没有被添加到陌生账户。
\end{enumerate}

\booksection{场景五：收到勒索信息——不要怕，不要付}

常见的勒索套路有两种：(1)“我有你的不雅视频/照片，不给钱就发给你通讯录里所有人”；(2)“你的手机被黑客控制了，不给钱就锁死”。遇到这种情况：

\begin{itemize}
  \item \textbf{99.9\%是诈骗}：他们手里根本没有你的任何隐私数据。群发同样的勒索信息给几万人，总有人害怕付钱。
  \item \textbf{不要回复，不要付钱}：回复了等于告诉对方"这个号码有人用"，你会被标记为"可勒索目标"，后续骚扰更多。
  \item \textbf{截图保存证据}：然后拉黑对方。
  \item \textbf{修改密码}：如果担心，修改一下你主要账号的密码，图个安心。
  \item \textbf{如果真的有隐私泄露风险}：报警。警方有处理这类案件的经验。
\end{itemize}

\booksection{场景六：AI换脸/AI拟声诈骗——眼见不再为实}

2024年起，AI换脸和AI拟声技术已经能让骗子在视频通话里“变成”你的家人、领导、朋友的脸和声音。“妈，我出事了，快转钱过来”——视频里的人长得、听起来都跟你儿子一模一样，但其实是假的。这是当前最难防、也最危险的诈骗手法之一。

\textbf{第一步：识别破绽——AI伪造的四个常见漏洞}\par

\begin{itemize}
  \item \textbf{眨眼频率异常}：真人每分钟眨眼15-20次，AI换脸往往眨眼过少或节奏僵硬。
  \item \textbf{嘴型与音频不同步}：仔细看口型和听到的声音之间，会有零点几秒的错位。
  \item \textbf{面部边缘模糊}：脸的轮廓、耳朵附近、头发边缘容易出现"毛边"或闪烁。
  \item \textbf{光影方向不一致}：脸上的高光、阴影和背景的光照方向对不上，或者侧脸时五官"跟着转"得不对。
\end{itemize}

\begin{tcbinfo}{第二步：三招验证对方真假}
\begin{itemize}
  \item \textbf{要求做特定动作}：让对方\textbf{转头、捂脸、用手挡住嘴巴}再做开合动作。AI换脸在遮挡和大角度转头时最容易"露馅"——脸会变形或跟不上动作。
  \item \textbf{问只有对方知道的问题}：编一个假问题，比如"上次咱家狗送去哪家医院看的？"——如果是骗子，他答不上来或者支支吾吾。
  \item \textbf{回拨已知号码}：挂断视频，用你通讯录里存的号码直接回拨过去。骗子控制不了真实的手机号，一拨就露馅。
\end{itemize}

\end{tcbinfo}

\begin{tcbwarn}{第三步：确认被骗后的报警流程}
\begin{enumerate}
  \item \textbf{立即停止转账}：再"急"也先停一下，凡是催你"快转账、别挂电话"的，几乎都是骗子。
  \item \textbf{保留全部证据}：录屏、截图保存视频通话画面、聊天记录、对方账号、转账凭证——这些是技术鉴定的关键材料。
  \item \textbf{拨打96110反诈专线}：告知是"AI换脸/拟声"类诈骗，接线员会启动紧急止付。
  \item \textbf{拨打110报警}：带齐证据去派出所做笔录，特别说明使用了AI伪造技术，便于警方立案定性。
  \item \textbf{通知亲友}：骗子可能用你的伪造视频去骗别人，提前在朋友圈或家庭群里发个提醒。
\end{enumerate}

\end{tcbwarn}

\booksection{场景七：SIM swap攻击（SIM劫持）——你的号码被“偷”走了}

这是一种很多人没听过、但危害极大的攻击。\textbf{SIM swap（SIM卡劫持）}指的是攻击者通过伪造你的身份信息，冒充你到营业厅补办一张新的SIM卡。一旦补办成功，你手机的信号会突然消失，而攻击者的手机插着那张新卡，能收到你所有的短信验证码——支付宝、银行、微信，全部门户大开。

\begin{tcbinfo}{怎么发现？SIM swap的三个预警信号}
\begin{itemize}
  \item 手机突然显示\textbf{"无服务"/"无SIM卡"}，但你明明没动过卡，重启也没用——这往往是被劫持的瞬间。
  \item 收到运营商发来的"补卡成功"或"业务办理成功"短信——但你本人没办过。
  \item 突然收不到验证码，但其他短信也收不到——因为短信被劫持到了攻击者的卡上。
\end{itemize}

\end{tcbinfo}

\begin{tcbwarn}{防护措施——把SIM卡也"上锁"}
\begin{itemize}
  \item \textbf{设置SIM卡PIN码}：路径\texttt{设置 > 电话 > SIM卡PIN码}。开启后，每次手机重启或换卡，都要输PIN才能用这张卡。攻击者就算拿到卡也用不了。
  \item \textbf{开启运营商二次验证}：登录中国移动/联通/电信App，找到"安全中心"或"业务办理验证"，开启\textbf{非本人补卡短信提醒}和\textbf{业务办理二次密码验证}。这样别人去补你的卡时，营业员会要求额外的身份核验。
  \item \textbf{重要账号别只用短信验证}：银行、Apple ID等高价值账号，优先用\textbf{认证器App}（Google Authenticator、微软Authenticator）或\textbf{硬件密钥}，而不是短信验证码。
  \item \textbf{身份证复印件加水印}：办理业务需要的复印件写明"仅供XX业务使用"——攻击者往往靠泄露的身份证信息冒充你。
\end{itemize}

\end{tcbwarn}

\textbf{被劫持后的应急处理——速度就是一切}\par

\begin{enumerate}
  \item \textbf{立即拨打运营商客服挂失}：移动10086 / 联通10010 / 电信10000，要求立即\textbf{停机并冻结补卡业务}，阻止攻击者继续使用你的号码。
  \item \textbf{用其他设备登录关键账号}：用电脑或家人手机，立即修改银行、支付宝、微信、Apple ID的密码，并踢掉陌生设备。
  \item \textbf{冻结支付账户}：支付宝95188挂失，微信安全中心冻结，银行客服冻结网银。
  \item \textbf{带身份证去营业厅}：本人持身份证到自营厅补办新卡，恢复号码使用权。
  \item \textbf{报警}：SIM swap已构成"侵犯公民个人信息罪"，带齐证据（停机时间记录、资金损失、补卡记录）报案。
  \item \textbf{全面排查}：恢复号码后，逐一检查近期登录记录和资金流水，确认有没有被冒用过的痕迹。
\end{enumerate}

\begin{tcbwarn}{预防胜于治疗——日常5件事}
\begin{enumerate}
  \item 锁屏密码 + 双重认证，两个都开
  \item "查找我的手机"确认开启
  \item SIM卡设置PIN码（设置 > 电话 > SIM卡PIN码），丢了手机别人不能用你的卡
  \item 不安装来历不明的App，不点陌生链接
  \item 定期备份重要数据（照片、通讯录、聊天记录）
\end{enumerate}

\end{tcbwarn}

\begin{tcbpractice}{动手：做一份属于你的"应急联系卡"}
\begin{enumerate}
  \item 找一张纸，写下以下号码：运营商客服（移动10086/联通10010/电信10000）、银行客服、96110反诈专线、110报警电话、家人电话。
  \item 写下你的Apple ID/品牌账号的注册邮箱和备用联系方式。
  \item 把这张纸放在钱包里或贴在冰箱上——不要存在手机里，因为出事时手机可能不在身边。
  \item 现在就检查："查找我的手机"是否已开启，SIM卡PIN码是否已设置。
\end{enumerate}

\end{tcbpractice}

\begin{tcbsummary}{本章小结}
\begin{itemize}
  \item 手机丢失黄金30分钟：借手机拨打远程锁定挂失SIM卡冻结支付通知亲友报警
  \item 被骗后72小时：停止转账保留证据拨96110报警联系银行改密码
  \item 手机中毒：安全模式卸载可疑App安全扫描恢复出厂（终极方案）
  \item 账号被盗：改密码踢设备查登录记录开双重认证查绑定查支付
  \item 勒索信息99.9\%是诈骗，不回复不付钱，截图拉黑
  \item AI换脸/拟声诈骗：看眨眼/嘴型/边缘/光影破绽，让对方转头捂脸、问私密问题、回拨已知号码
  \item SIM swap劫持：手机突然"无服务"即警觉；设SIM卡PIN、开运营商二次验证、重要账号改用认证器App
  \item 日常预防五件事：锁屏+2FA、查找手机、SIM卡PIN、不装来历不明App、定期备份
\end{itemize}

\end{tcbsummary}
